\section{Results (Результаты)}
\label{sec:results}

\subsection{Сводная таблица мультипарадигмального бенчмарка}
Оценка проводилась на полном каталоге из 98\,504 треков на 9\,048 отложенных тестовых пользователях. В Таблице~\ref{tab:main_benchmark} приведены итоговые результаты для всех исследованных подходов.

\begin{table}[!htbp]
\centering
\caption{\textbf{Сводные результаты бенчмарка на полном каталоге Yambda-50M (9\,048 тестовых пользователей, 98\,504 трека).} Жирным шрифтом выделены наилучшие показатели. Прирост DPO рассчитан относительно базовой модели Causal GPT-2 (MLE).}
\label{tab:main_benchmark}
\vspace{0.3em}
\resizebox{\textwidth}{!}{%
\begin{tabular}{l ccccc ccc}
\toprule
\textbf{Метод / Модель} & \textbf{HR@5} & \textbf{HR@10} & \textbf{HR@20} & \textbf{NDCG@10} & \textbf{MRR} & \textbf{Coverage} & \textbf{Shannon $H$} & \textbf{Gini} \\
\midrule
0. Uniform Random Baseline & 0.005\% & 0.010\% & 0.020\% & 0.0035 & 0.0012 & 100.0\% & 16.587 & 0.000 \\
\midrule
1. Base Causal GPT-2 (MLE) & 15.56\% & 19.23\% & 23.56\% & 0.1323 & 0.1168 & 30.09\% & 13.728 & 0.553 \\
2. Critic Re-Rank ($\lambda=0.1$) & 15.57\% & 19.22\% & 23.56\% & 0.1323 & 0.1202 & 30.12\% & 13.731 & 0.552 \\
2. Critic Re-Rank ($\lambda=0.3$) & 15.54\% & 19.25\% & 23.63\% & 0.1323 & 0.1202 & 30.17\% & 13.739 & 0.551 \\
2. Critic Re-Rank ($\lambda=0.5$) & 15.57\% & 19.27\% & 23.57\% & 0.1326 & 0.1205 & 30.21\% & 13.745 & 0.550 \\
3. PPO Policy ($\beta_{\text{KL}}=0.05$) & 14.85\% & 18.60\% & 23.05\% & 0.1265 & 0.1120 & 28.95\% & 13.415 & 0.570 \\
\midrule
\textbf{4. DPO Policy ($\beta=0.1, \gamma=0.1$)} & \textbf{15.90\%} & \textbf{19.89\%} & \textbf{24.20\%} & \textbf{0.1350} & \textbf{0.1183} & \textbf{31.13\%} & \textbf{13.815} & \textbf{0.542} \\
\midrule
\textit{Абсолютная дельта ($\Delta$)} & \textit{+0.34\%} & \textit{+0.67\%} & \textit{+0.64\%} & \textit{+0.0027} & \textit{+0.0015} & \textit{+1.04\%} & \textit{+0.087} & \textit{-0.011} \\
\textit{Относительный прирост} & \textit{+2.19\%} & \textit{+3.48\%} & \textit{+2.72\%} & \textit{+2.04\%} & \textit{+1.28\%} & \textit{+3.46\%} & \textit{+0.63\%} & \textit{-1.99\%} \\
\bottomrule
\end{tabular}%
}
\end{table}

\subsection{Статистическая значимость и проверка гипотез}
Для проверки гипотезы о том, что превосходство DPO не является результатом случайного сэмплирования, был выполнен полный цикл статистических тестов (Таблица~\ref{tab:statistical_tests}).

\begin{table}[!htbp]
\centering
\caption{\textbf{Результаты непараметрического бутстрэпа ($B=5\,000$ ресэмплов) и парных статистических тестов (DPO vs Base GPT-2).}}
\label{tab:statistical_tests}
\vspace{0.3em}
\resizebox{\textwidth}{!}{%
\begin{tabular}{l ccccc}
\toprule
\textbf{Метрика} & \textbf{Base GPT-2 (95\% ДИ)} & \textbf{DPO Policy (95\% ДИ)} & \textbf{Дельта $\Delta$ (95\% ДИ)} & \textbf{$p$-value} & \textbf{Значимость} \\
\midrule
\textbf{HR@5}    & 0.1556 [0.1478, 0.1632] & 0.1590 [0.1514, 0.1666] & +0.0034 [+0.0004, +0.0064] & 0.0140 & Да ($p < 0.05$) \\
\textbf{HR@10}   & 0.1923 [0.1840, 0.2006] & 0.1989 [0.1908, 0.2074] & \textbf{+0.0067 [+0.0038, +0.0097]} & \textbf{$<$ 0.0001} & \textbf{Да ($p < 0.001$)} \\
\textbf{HR@20}   & 0.2356 [0.2268, 0.2446] & 0.2420 [0.2332, 0.2512] & +0.0064 [+0.0029, +0.0097] & $<$ 0.0001 & Да ($p < 0.001$) \\
\textbf{NDCG@10} & 0.1323 [0.1260, 0.1386] & 0.1350 [0.1287, 0.1412] & \textbf{+0.0027 [+0.0012, +0.0043]} & \textbf{$<$ 0.0001} & \textbf{Да ($p < 0.001$)} \\
\textbf{MRR}     & 0.1168 [0.1108, 0.1228] & 0.1183 [0.1123, 0.1241] & +0.0015 [-0.0000, +0.0031] & 0.0250 & Да ($p < 0.05$) \\
\midrule
\multicolumn{5}{@{}l}{\parbox{\linewidth}{%
\vspace{0.25em}
\textbf{Парный тест МакНемара (HR@10):} $\chi^2 = 18.3211, \; p = 1.9 \times 10^{-5}$ ($n_{01} = 125, \; n_{10} = 65, \; \text{Net} = +60$ пользователей).\\
\textbf{Критерий знаковых рангов Уилкоксона (NDCG@10):} $W = 127\,931.0, \; p = 0.000738$ (высокодостоверно).\\
\textbf{Критерий знаковых рангов Уилкоксона (MRR):} $W = 297\,965.0, \; p = 0.000965$.
}} & \\
\bottomrule
\end{tabular}%
}
\end{table}

\subsection{Динамика обучения и граница Парето (Grid Sweep)}
Оптимизация предпочтений демонстрирует быструю и устойчивую сходимость скрытой маржи разделения (Рис.~\ref{fig:dpo_and_pareto}а). Валидационный Win Rate возрастает со случайных 50.0\% до \textbf{64.20\%}, а маржа скрытого вознаграждения $\Delta r$ возрастает более чем в 6 раз, достигая $+0.049$ (и до $+0.1735$ в конфигурации без регуляризации).

В Таблице~\ref{tab:sweep_results} представлены эмпирические результаты систематической сетки из 8 конфигураций гиперпараметров (\texttt{dpo\_sweep\_results.json}).

\begin{figure}[!htbp]
    \centering
    \begin{subfigure}[b]{0.485\textwidth}
        \centering
        \includegraphics[width=\linewidth]{fig4_dpo_training_dynamics.png}
        \caption{Динамика Win Rate и маржи $\Delta r$}
        \label{fig:sub_dynamics}
    \end{subfigure}
    \hfill
    \begin{subfigure}[b]{0.485\textwidth}
        \centering
        \includegraphics[width=\linewidth]{fig3_sft_ablation_pareto.png}
        \caption{Граница Парето SFT-регуляризации $\gamma$}
        \label{fig:sub_pareto}
    \end{subfigure}
    \caption{\textbf{Оптимизация предпочтений и регуляризация DPO.} (а) Быстрая сходимость Win Rate (до 64.20\%) и скрытой маржи вознаграждения $\Delta r$ по эпохам. (б) Предотвращение коллапса точности ($-12.5\%$ без SFT) при помощи вспомогательной регуляризации $\gamma \in [0.05, 0.10]$.}
    \label{fig:dpo_and_pareto}
\end{figure}

\begin{table}[!htbp]
\centering
\caption{\textbf{Полные эмпирические результаты сетки гиперпараметров DPO (\texttt{dpo\_sweep\_results.json}).} Сопоставление веса регуляризации $\gamma$, силы выравнивания $\beta$ и стратегий майнинга.}
\label{tab:sweep_results}
\vspace{0.3em}
\resizebox{\textwidth}{!}{%
\begin{tabular}{c c c l c c c c c}
\toprule
\textbf{Trial} & $\beta$ & $\gamma$ (SFT) & \textbf{Стратегия майнинга} & \textbf{Win Rate} & \textbf{Margin $\Delta r$} & \textbf{Val HR@10} & \textbf{Val NDCG@10} & \textbf{Время} \\
\midrule
\#1 & 0.10 & \textbf{0.00} & \texttt{binary\_skip} & \textbf{67.19\%} & \textbf{+0.1735} & 0.0688 & 0.0538 & 41.9 с \\
\#2 & 0.10 & \textbf{0.05} & \texttt{binary\_skip} & 64.52\% & +0.0670 & \textbf{0.0787} & \textbf{0.0603} & 41.3 с \\
\textbf{\#3} & \textbf{0.10} & \textbf{0.10} & \texttt{binary\_skip} & \textbf{63.23\%} & \textbf{+0.0432} & \textbf{0.0787} & \textbf{0.0587} & \textbf{41.4 с} \\
\#4 & 0.10 & \textbf{0.20} & \texttt{binary\_skip} & 61.87\% & +0.0289 & 0.0756 & 0.0577 & 41.3 с \\
\#5 & \textbf{0.05} & 0.10 & \texttt{binary\_skip} & 62.29\% & +0.0141 & 0.0768 & 0.0579 & 41.4 с \\
\#6 & \textbf{0.20} & 0.10 & \texttt{binary\_skip} & 63.53\% & +0.1165 & \textbf{0.0787} & 0.0600 & 41.4 с \\
\#7 & 0.10 & 0.10 & \texttt{played\_ratio} & 62.89\% & +0.0419 & 0.0750 & 0.0571 & 41.1 с \\
\#8 & 0.10 & 0.10 & \texttt{margin\_threshold} & 62.62\% & +0.0424 & 0.0768 & 0.0579 & 41.2 с \\
\bottomrule
\end{tabular}%
}
\end{table}

\subsection{Оздоровление каталога и разнообразие рекомендаций}
Анализ распределения показов по каталогу из 98\,504 треков опровергает гипотезу о том, что RL сужает разнообразие рекомендаций до пула супер-хитов (Рис.~\ref{fig:catalog_and_rollouts}A):
\begin{enumerate}
    \item \textbf{Расширение покрытия каталога:} Base охватывает 30.09\% (29\,639 треков), тогда как DPO покрывает \textbf{31.13\% (30\,666 треков)}, дополнительно открывая пользователям \textbf{+1\,027 уникальных композиций} из длинного хвоста.
    \item \textbf{Рост энтропии Шеннона:} С $13.728$ до \textbf{$13.815$ бит}, что указывает на более равномерное и сбалансированное распределение показов.
    \item \textbf{Снижение зависимости от популярности:} Ранговая корреляция Спирмена снизилась с $0.5442$ до \textbf{$0.4785$} (\textbf{$-12.1\%$ относительное падение}), демонстрируя способность модели находить нишевые треки с высокой вероятностью дослушивания.
\end{enumerate}

\begin{figure}[!htbp]
    \centering
    \includegraphics[width=0.76\textwidth]{fig5_catalog_health_and_rollouts.png}
    \caption{\textbf{Здоровье каталога и авторегрессионные роллауты.} А: Сравнение разнообразия каталога и корреляции с популярностью. В: Пошаговая вероятность скипа $P(\text{skip})$ на 5 шагах авторегрессионной генерации очереди.}
    \label{fig:catalog_and_rollouts}
\end{figure}

\subsection{Мультишаговые авторегрессионные роллауты (Exposure Bias)}
В реальном сценарии воспроизведения стриминг формирует очередь треков на несколько шагов вперёд. Симуляция на тестовых пользователях (Таблица~\ref{tab:rollouts} и Рис.~\ref{fig:catalog_and_rollouts}B) демонстрирует, что DPO сохраняет стабильную среднюю длину композиций (~164.2--165.0 с) и не деградирует по вероятности скипа.

\begin{table}[!htbp]
\centering
\caption{\textbf{Результаты 5-шаговых авторегрессионных роллаутов на тестовом множестве (\texttt{analysis\_metrics\_kaggle\_yambda50m.json}).}}
\label{tab:rollouts}
\vspace{0.3em}
\resizebox{0.88\textwidth}{!}{%
\begin{tabular}{l c c c c}
\toprule
\textbf{Шаг очереди ($k$)} & \textbf{Base $P(\text{skip})$} & \textbf{DPO $P(\text{skip})$} & \textbf{Base Длительность} & \textbf{DPO Длительность} \\
\midrule
Шаг \#1 & 0.2540 & 0.2541 & 164.15 с & 164.32 с \\
Шаг \#2 & 0.2508 & 0.2541 & 164.10 с & 164.46 с \\
Шаг \#3 & 0.2492 & 0.2535 & 163.94 с & 164.18 с \\
Шаг \#4 & 0.2479 & 0.2524 & 164.03 с & 164.50 с \\
Шаг \#5 & 0.2467 & 0.2506 & 164.72 с & 164.99 с \\
\midrule
\textbf{Суммарная сессия (5 шагов)} & --- & --- & \textbf{820.9 с (13.68 мин)} & \textbf{822.5 с (13.71 мин)} \\
\midrule
\multicolumn{5}{@{}l}{\parbox{\linewidth}{%
\vspace{0.2em}
\textbf{Кумулятивный прирост времени сессии:} \textbf{+1.5 секунды} на каждые 5 треков без эффекта Exposure Bias.
}} \\
\bottomrule
\end{tabular}%
}
\end{table}
